\section{Evaluation design}
\label{sec:experiments}

\paragraph{Pretraining and held-out evaluation.}
ZoomST is pretrained on 45 MOSTA mouse embryo sections~\citep{chen2022mosta}, with eight sections reserved in advance, one from each sampled stage from E9.5 to E16.5. Supervision uses only granularities 8, 32, and 128. Held-out representation and transformation tests freeze the encoders and transition; input preprocessing and target-fitted readouts follow Appendices~\ref{app:preprocessing} and~\ref{app:dual-view}. The No Cross ablation omits cross-granularity supervision with the same total number of training updates. Held-out tests evaluate generalization, matched continuations analyze the training objectives, and source-section anatomical cases interpret granularity-dependent representations and responses. Table~\ref{tab:evaluation_protocol} specifies data use for each analysis.

\paragraph{Retained tissue information and correspondence.}
Whole-section clustering compares frozen ZoomST, Novae, and SToFM with target-fitted GraphST, BANKSY, MENDER, and STAGATE, sharing positions, annotation-defined cluster counts, and downstream readouts with the stated method-specific preprocessing. AMI and physical 32-neighbor retention measure anatomical and local spatial organization. Stage-balanced Brain observations from the held-out sections evaluate developmental ordering by diffusion-pseudotime--stage correlations (Appendix~\ref{app:developmental-ordering}). Separately, equal-update continuations on three sections compare Base, Base+CS, and Base+CS+JointMMD on their respective adapted sections. Native neighbor overlap and same-location Recall@32 measure molecular-state correspondence without a fitted transformation (Appendix~\ref{app:objective-ablation}).

\paragraph{Transformation and response generalization.}
Direct neighborhood encodings provide reference targets. Trained-pair state prediction uses all held-out positions; joint-association tests include shuffled-pair controls and projections unused in training. Unseen-target prediction uses 1,024 fixed positions per section at granularities 12, 16, 24, 48, 64, and 96. These target encodings do not enter training or prediction. The same held-out locations evaluate response ordering over four doubling intervals (Appendices~\ref{app:granularity-transfer} and~\ref{app:scale-continuum}).

\paragraph{Anatomical interpretation.}
Matched IoU readouts compare Lung and Smooth muscle at fine and coarse granularities. Facial nerve and Brain response analyses compare predicted change with encoded change and annotated neighborhood composition. These anatomical cases use E14.5 E1S1 and E16.5 E1S1 from the training set; the Brain mapping display uses held-out E16.5 E2S13. Physical extents use the exact nested supports and nominal bin50 calibration. Sampling, domains, and response definitions are specified in Appendices~\ref{app:scale-cases} and~\ref{app:response-meaning}.
